\section{Final disposition of the treated questions}

No unchecked candidate is used as a conclusion of this paper. Each row of
\texttt{QUESTIONS.tsv} is resolved by a stated theorem,
counterexample, or correction of the precise inference asked about.
The explicit input qualification in R45 is part of its theorem: a numerical
table for an unnamed member and an unnamed rational Weil vector is not
mathematical data supplied by T5. P12 and V04 prove why the prime and Newton
type alone cannot determine such a table. R45 supplies its exact terminating
calculation once those data are specified.

The earlier working interval $16\leq c_{\C,\theta}(w)\leq24$ is closed
by W04. Its proof is computer-assisted and concerns the complex
normalization stated there. B03 also gives its eventual local value
twenty-four. The all-good-prime values sixteen for $t_*$ and $t_0$ are
separate, fully geometric theorems P13 and V06. No assertion about an
unevaluated exceptional good prime for the third normalization is inserted
into either all-prime theorem.

The attempted unrestricted finite-characteristic period inequality is
refuted by P14. The graph-only catalog attempted for the final
degree-twenty-two comparison is corrected by W03 and is replaced by the
complete W04 catalog. The same-family CM-transport inference in the
sketch is rejected by C02 as stronger than its source theorem; no
general nonexistence theorem for such CM points is claimed. A
prime-growth law for arbitrarily chosen transport data is refuted by
C06, while C03--C05 give the exact charged data and C18 supplies the
requested explicit constant-growth example.

\section{Historical identifiers and accepted proof locations}

The research memos preserve earlier identifiers so that their actual
checks can be traced. They do not add unstated final theorems. The
following identifications avoid counting an independently checked step
twice as a new result:
\begin{itemize}
\item C08 is the algebraic coordinate construction preceding P13;
C09--C10 are the low-degree geometric lemma and its curve-image
factorization there. C11--C12 are its characteristic-two and
characteristic-three calculations. C13 is the all-good-prime theorem
P13, and C14 is the counterexample P14. Their full proofs are included
in the local-Weil section.
\item F01--F10 are the independent re-proof of the low-degree
components of P13. F11 records the scope of a corroborating finite
enumeration; the final P13 proof does not depend on that enumeration.
\item P01 and P04 record checks of the cited Hodge and arithmetic
inputs. P05 is the fixed-algebraic-class specialization consequence of
B03. P10's original bounded interval is superseded by W04. P11's
distinction between prime distribution and a cycle-degree bound is
made precise in R30--R31.
\item The suffixes in the R20--R24 proof-check memos refer to the
corresponding numbered steps in the final sixfold-extension section.
C33 and its suffixes are proof-check events for V01--V06, not new
theorems.
\end{itemize}

The unfinished Hecke statement indexed in the supplied Masullo material
is recorded with its actual conditional status and is not used. The
paper makes no claim to solve the higher-dimensional Hecke-orbit problems
assigned to the parallel research runs. No findings from those runs
were supplied for verification.

All mathematical questions treated here reached the stated dispositions
before the first scheduled three-hour progress review. The final ledger
therefore records completion, without inventing a timed stagnation review
or invoking the two-review stopping rule.
